\chapter{Anexos}

\section{Anexo A. Referencia de endpoints del sistema}
La tabla~\ref{tab:apiref} recoge los endpoints transversales del framework
(no ligados a una entidad concreta) que conforman la superficie de
descubrimiento y esquema del contrato.

\begin{table}[!htbp]
\caption{Endpoints transversales del contrato.}
\label{tab:apiref}
\centering
\begin{tabular}{|l|l|p{6cm}|}
\hline
\textbf{Método} & \textbf{Ruta} & \textbf{Propósito} \\ \hline
GET & \texttt{/api/sys/entities} & Entidades CRUD y sus capacidades \\ \hline
GET & \texttt{/api/sys/discovery} & Servicios externos (Traefik) \\ \hline
GET & \texttt{/api/sys/schemas} & Bundle versionado de esquemas (ETag) \\ \hline
GET & \texttt{/<entity>/schema.json} & Esquema enriquecido de la entidad
(ETag, 304) \\ \hline
PUT & \texttt{/api/authn?user=...} & Login local de desarrollo \\ \hline
GET & \texttt{/health} & Sonda de estado del servicio \\ \hline
\end{tabular}
\end{table}

\section{Anexo B. Referencia de configuración}
La inicialización acepta las secciones de configuración resumidas en la
tabla~\ref{tab:configref}, validadas por Pydantic Settings y sobreescribibles
mediante variables de entorno con prefijo \texttt{UNIBACK\_}.

\begin{table}[!htbp]
\caption{Secciones de configuración del inicializador.}
\label{tab:configref}
\centering
\begin{tabular}{|l|p{8cm}|}
\hline
\textbf{Sección} & \textbf{Claves principales} \\ \hline
\texttt{database} & \texttt{url}, \texttt{echo}, \texttt{pool\_size} \\ \hline
\texttt{api} & \texttt{title}, \texttt{version}, \texttt{debug},
\texttt{cors\_origins} \\ \hline
\texttt{auth} & \texttt{secret\_key}, \texttt{algorithm},
\texttt{firebase\_credentials\_path} \\ \hline
\texttt{worker} & \texttt{enabled}, \texttt{broker\_url},
\texttt{backend\_url} \\ \hline
\end{tabular}
\end{table}

\section{Anexo C. Acceso al código fuente}
El código completo del framework reside en dos repositorios: el backend
\emph{uniback} (Python, paquete instalable con \texttt{pip install -e
".[dev]"}) y el frontend \emph{gui-components}, que contiene la librería
publicable \texttt{ngt-gui}, en la que el componente
\texttt{<ngt-dynamic-form>} integra el contrato extremo a extremo. La
documentación interna
del contrato y de los incrementos se encuentra en la carpeta \texttt{docs/}
del backend (\texttt{CONTRACT.md}, \texttt{DYNAMIC\_FORMS.md},
\texttt{SCHEMA\_GENERATION.md}), y constituye la guía de referencia para los
desarrolladores que reutilicen el framework.

\section{Anexo D. Demo y herramientas añadidas}
Este anexo documenta dos incrementos posteriores a la consolidación del
núcleo, ambos demostrables sobre el despliegue multinodo y realizados sin
alterar el contrato REST ni el frontend existentes: la migración del backend a
una arquitectura de micronúcleo y la incorporación de un asistente de
inteligencia artificial multimodelo.

\subsection{D.1. Migración a una arquitectura de micronúcleo}
Partiendo de un paquete \emph{uniback} internamente monolítico (la factoría de
la aplicación montaba una treintena de \emph{routers} mediante condicionales,
el inicializador importaba todos los modelos y el sembrado mezclaba todos los
dominios), se refactorizó el backend hacia un \textbf{micronúcleo}: el paquete
queda reducido al núcleo y cada dominio funcional pasa a ser un plugin de serie
en \texttt{uniback.contrib.<dominio>}, conservando todas las capacidades.

El núcleo retiene la persistencia (modelos \texttt{core}, \texttt{sysadmin} y
\texttt{screens}), la autenticación y autorización (la dependencia
\texttt{get\_n\_session} y el control de acceso por ACL), la factoría CRUD, el
registro de esquemas (\texttt{schema\_registry}), el sembrado del propio núcleo
y los servicios de sistema (sonda de salud, navegación de la interfaz,
descubrimiento, importación genérica, objetos funcionales y funciones de
sistema). Cada dominio --- autenticación, ficheros, anotaciones, especies,
jerarquías, colecciones, CRUD de interfaz y geografía --- aporta sus propios
modelos, \emph{routers} y sembrado como plugin.

La activación por nodo es la pieza central: cada plugin declara los nodos en
los que está activo (atributo \texttt{node\_types}) y la variable de entorno
\texttt{UNIBACK\_NODE\_TYPE} decide qué plugins montan sus \emph{routers} y
siembran (el valor \texttt{monolith} los activa todos). El matiz clave es que
\emph{todos} los nodos cargan \emph{todos} los modelos --- requisito del
polimorfismo de \texttt{FunctionalObject} --- y registran el catálogo completo
de entidades en el \texttt{schema\_registry} (para \texttt{/api/sys/schemas} y
las pantallas sintéticas), pero solo \emph{montan} los \emph{routers} de los
plugins activos; Traefik enruta cada ruta al nodo que la sirve. La
correspondencia entre dominios y nodos se resume en la
tabla~\ref{tab:dominios-nodo}.

\begin{table}[!htbp]
\caption{Distribución de dominios funcionales por tipo de nodo.}
\label{tab:dominios-nodo}
\centering
\begin{tabular}{|l|l|}
\hline
\textbf{Dominio (plugin)} & \textbf{Nodo} \\ \hline
\texttt{auth} & \texttt{auth} \\ \hline
\texttt{files} & \texttt{files} \\ \hline
\texttt{annotations} & \texttt{annotations} \\ \hline
\texttt{species} & \texttt{species} \\ \hline
\texttt{hierarchies}, \texttt{collections}, & \multirow{2}{*}{\texttt{core}} \\
\texttt{gui\_crud}, \texttt{geographics} & \\ \hline
\texttt{assistant} & \texttt{assistant} \\ \hline
\end{tabular}
\end{table}

La compatibilidad se preserva mediante redirecciones de retrocompatibilidad
(\emph{shims}) en \texttt{persistence/models}, de modo que no cambian ni la API
REST, ni las etiquetas de enrutado de Traefik, ni el frontend.

Durante las pruebas se detectó y corrigió un problema de aislamiento
relacionado: un plugin de tipo \emph{hot-plug} debe declarar su
\texttt{node\_types}; en caso contrario su \emph{router} se monta en todos los
nodos y, al detener su nodo dedicado, Traefik hace caer la ruta en el comodín
del nodo \texttt{core}, que la seguiría sirviendo. La entidad quedaría así
accesible por URL pese a estar «desenchufada». Restringiendo el plugin a su
propio nodo, el endpoint solo lo sirve dicho nodo (devolviendo \texttt{404}
cuando está caído), mientras la entidad permanece en el catálogo de esquemas.

\subsection{D.2. Asistente de IA multimodelo}
\label{anexo:ia}
Como segunda capacidad se añadió un asistente de inteligencia artificial, de
nuevo como plugin de serie (\texttt{uniback.contrib.assistant}, nodo
\texttt{assistant}), sin tocar el núcleo. En el frontend se presenta como una
ventana superpuesta, minimizable y arrastrable que se ancla a la esquina más
cercana. El asistente interactúa con la aplicación: consulta datos, dirige al
usuario a la pantalla donde se encuentran y propone altas en la tabla que está
visualizando.

Para que \textbf{varios modelos coexistan y sean seleccionables} se
introdujeron dos registros extensibles, gemelos de los ya empleados por el
framework para sus extensiones: uno de proveedores de modelo
(\texttt{model\_provider\_registry}) y otro de herramientas
(\texttt{assistant\_tool\_registry}). De serie se registran los modelos Claude
(Opus~4.8 por defecto, Sonnet~4.6 y Haiku~4.5) mediante el SDK oficial
\texttt{anthropic}, y se admiten proveedores propios \emph{OpenAI-compatible}
configurables por URL y clave de API para una IA local o de empresa
--- verificado con \emph{LM Studio} ---, según el ejemplo del
listado~\ref{lst:proveedor-local}. El frontend descubre y selecciona de entre
lo integrado en el backend mediante los endpoints
\texttt{GET /api/assistant/models} y \texttt{GET /api/assistant/tools}.

\begin{lstlisting}[caption={Proveedor de modelo local configurado por variable de entorno.}, label={lst:proveedor-local}]
[{"name": "lmstudio-gemma",
  "label": "Gemma (LM Studio)",
  "base_url": "http://host.docker.internal:1234/v1",
  "api_key": "lm-studio",
  "model_id": "google/gemma-4-e4b"}]
\end{lstlisting}

Las herramientas son de dos clases. Las de \emph{servidor}
(\texttt{list\_entities}, \texttt{get\_entity\_schema} y
\texttt{query\_records}) se ejecutan en el backend bajo la sesión del usuario
y, por tanto, respetan el ACL del núcleo sin necesidad de lógica adicional. Las
de \emph{interfaz} (\texttt{navigate\_to} y \texttt{propose\_table\_rows}) las
ejecuta el \emph{overlay} en el navegador. El endpoint de conversación
(tabla~\ref{tab:asistente-endpoints}) es un flujo de eventos en \emph{streaming}
(SSE) protegido por \texttt{get\_n\_session}; un bucle de uso de herramientas
agnóstico del proveedor ejecuta las de servidor y emite al cliente las de
interfaz.

\begin{table}[!htbp]
\caption{Endpoints del asistente de IA.}
\label{tab:asistente-endpoints}
\centering
\begin{tabular}{|l|l|p{6cm}|}
\hline
\textbf{Método} & \textbf{Ruta} & \textbf{Propósito} \\ \hline
GET & \texttt{/api/assistant/models} & Modelos seleccionables (del registro) \\ \hline
GET & \texttt{/api/assistant/tools} & Herramientas integradas (del registro) \\ \hline
POST & \texttt{/api/assistant/chat} & Chat con \emph{streaming} (SSE), bajo la sesión del usuario \\ \hline
\end{tabular}
\end{table}

En cuanto a la seguridad, el control de acceso se hereda al ejecutar bajo la
sesión del usuario que conversa; las escrituras siguen un patrón
\textbf{proponer-y-confirmar} --- el modelo propone las filas y el alta real
solo se efectúa, a través del CRUD existente, cuando el usuario confirma ---; y
el contenido de los registros se trata como datos, nunca como instrucciones.
